% ── 5.2  H1: Resolution and the Performance--Cost / Footprint Frontiers ──────
% PURPOSE: The signature analysis -- whether forgetting matches full memory on the
%          CL-F1 proxy, and how the two cost axes (compute vs footprint) behave.
% INTEGRITY: every numeric value and every conclusion-clause is a PLACEHOLDER pending
% the trustworthy 144 rerun (runs_144_seq_cceb325). The pre-registered ANALYSIS PLAN is
% writable now; the ANSWER is not. Primary test = Wilcoxon (Invariant #11); TOST/SESOI
% was NOT computed -- non-inferiority is read from bounded CIs + power, not a TOST verdict.
% Figures are placeholder boxes (results/plots is regenerated from the rerun).
% ─────────────────────────────────────────────────────────────────────────

\section{H1: Resolution and the Performance--Cost / Footprint Frontiers}
\label{sec:h1}

This section reports the thesis's primary contrast: whether any forgetting policy differs
from full-memory accumulation on the resolution proxy, and how cost behaves on the compute
and footprint axes. Because retrieval is held constant across all six conditions, the
contrast is attributable to \emph{what} each policy keeps rather than to \emph{how} memory
is searched. The analysis plan below is fixed in advance; the quantitative verdict is
filled from the trustworthy rerun \texttt{[RESULT: pending runs\_144\_seq\_cceb325]}.

\subsection{The primary test}

The pre-registered primary analysis is a Wilcoxon signed-rank test on $N=8$
sequence-level means, comparing each of the five other policies against
\texttt{full\_memory}, with Holm correction across the five contrasts
(Section~\ref{sec:metrics}). Effect size is rank-biserial $r_{rb}$ and the interval is the
95\% BCa bootstrap CI on the paired difference in resolution rate.
Table~\ref{tab:main-results} reports the five contrasts once the rerun completes; the
direction and significance of each are \texttt{[RESULT: pending]}.

\begin{table}[htbp]
  \centering
  \caption{H1 primary test. Wilcoxon signed-rank on $N=8$ sequence means, each policy vs
  \texttt{full\_memory}, Holm-corrected over the five contrasts. Effect size is
  rank-biserial $r_{rb}$; the interval is the 95\% BCa bootstrap CI on the paired
  difference in resolution rate. \texttt{[RESULT: pending runs\_144\_seq\_cceb325]}}
  \label{tab:main-results}
  \begin{tabular}{lrrrrlc}
    \toprule
    Contrast (vs \texttt{full\_memory}) & $W$ & $p$ & Holm $p$ & $r_{rb}$ & 95\% BCa CI & Sig. \\
    \midrule
    random\_prune      & --- & --- & --- & --- & --- & --- \\
    recency\_prune     & --- & --- & --- & --- & --- & --- \\
    type\_aware\_decay & --- & --- & --- & --- & --- & --- \\
    cls\_consolidation & --- & --- & --- & --- & --- & --- \\
    no\_memory         & --- & --- & --- & --- & --- & --- \\
    \bottomrule
  \end{tabular}
\end{table}

\subsection{Bounding and powering the result}

A non-significant result is informative only if the test could have detected an effect
worth caring about; a significant one needs its magnitude bounded. We therefore report,
alongside the primary test, the aggregate full-memory-minus-no-memory difference with its
bootstrap interval and the detectable-effect (power) profile of the $N=8$ Wilcoxon design
under the observed sequence-level noise. These are presented in the estimation spirit of
\cite{cumming2014} and \cite{wasserstein2019} --- as intervals and detectable-effect
sizes that bound what the null does or does not rule out, \emph{not} as a second
significance verdict, and with no test-switching to manufacture significance (Invariant
\#11). The magnitudes (the bounded memory benefit, the per-percentage-point power, and the
per-task memory-churn versus seed-noise comparison) are
\texttt{[RESULT: pending runs\_144\_seq\_cceb325]}.

\subsection{The two-axis frontier: compute versus footprint}

The efficiency contribution (H1b) is read on two separate axes, because forgetting can
behave very differently on each.

\paragraph{Compute.}
On total token cost (Figure~\ref{fig:pareto-compute}) the question is whether the six
policies separate at all, given that cost is dominated by the agent loop --- each task
runs up to 20 ReAct steps whose growing conversation is resent every step, so the memory
payload is a small fraction of the per-task token bill. The per-policy spread relative to
the paired standard deviation is \texttt{[RESULT: pending runs\_144\_seq\_cceb325]}.

\paragraph{Footprint.}
The memory-footprint axis (Figure~\ref{fig:pareto-footprint}) measures the memory tokens
carried per task --- the quantity that governs storage and retrieval-index growth over a
long-running agent's lifetime, and the axis on which the pruning policies are designed to
differ from \texttt{full\_memory}. The footprint ratio between full accumulation and the
pruners, and whether any policy is Pareto-dominated, are
\texttt{[RESULT: pending runs\_144\_seq\_cceb325]}.

\begin{figure}[htbp]
\centering
\begin{figurestub}{Performance vs compute (pending)}
Two-axis Pareto: CL-F1 (resolution rate) against total tokens per run, with 95\%
confidence ellipses per policy. Regenerated from \texttt{results/plots/pareto\_compute.png}
once the trustworthy rerun (\texttt{runs\_144\_seq\_cceb325}) is aggregated.
\end{figurestub}
\caption{Performance versus compute (total tokens per run).
\texttt{[RESULT: pending runs\_144\_seq\_cceb325]}}
\label{fig:pareto-compute}
\end{figure}

\begin{figure}[htbp]
\centering
\begin{figurestub}{Performance vs memory footprint (pending)}
Two-axis Pareto: CL-F1 against memory tokens carried per task. \texttt{full\_memory} sits
at the high-footprint end; the pruning policies and \texttt{no\_memory} populate the
light-footprint corner. Regenerated from
\texttt{results/plots/pareto\_footprint.png} once the rerun is aggregated.
\end{figurestub}
\caption{Performance versus memory footprint (tokens carried per task).
\texttt{[RESULT: pending runs\_144\_seq\_cceb325]}}
\label{fig:pareto-footprint}
\end{figure}

\paragraph{The combined claim.}
Read together, the value axis and the two cost axes determine whether aggressive
forgetting matches full-memory accumulation on task resolution while carrying a smaller
footprint, and at what compute cost. The synthesised statement is deferred until the
numbers land \texttt{[RESULT: pending runs\_144\_seq\_cceb325]}; the rest of the chapter
quantifies \emph{why} the value-axis gap is the size it is.
